\chapter{回転、スピノル、色対称性}

\section{群としての回転}

空間の回転は対称性であり、さまざまな種類の対象に作用する。剛体、速度ベクトル、ベクトル場、量子状態に作用する。その変換を直接思い描けるので、通常のベクトルから始めるのが最も簡単である。

回転は組み合わせることもできる。ある軸のまわりに或る角度だけ物体を回転させ、次に第二の軸のまわりに第二の角度だけ回転させる。最終的な向きは、或る第三の軸のまわりに或る第三の角度だけ回転させる単一の回転と等価である。したがって、二つの回転の積はまた一つの回転である。

この積は変換に対する演算であり、二つの軸ベクトルの和ではない。第一の回転は、第二の回転が作用する対象を変える。\(R_1\)を先に施し、\(R_2\)を次に施すなら、合成された演算は\(R_2R_1\)と書かれる。それはそれ自体が或る\(R_3\)である。ただし\(R_3\)の軸と角度を求めるのは厄介かもしれない。合成の下でのこの閉性が、回転を単なる幾何学的運動の一覧ではなく群とみなす第一の理由である。

まず一つの回転を幾何学的に記述すると有用である。回転軸に沿った単位ベクトル\(\hat{\boldsymbol n}\)と角度\(\theta\)を選ぶ。右手の規約を採用する。\(\theta\)だけの左手系の回転は、\(-\theta\)だけの右手系の回転として、あるいは\(-\hat{\boldsymbol n}\)のまわりの右手系の回転として等しく書くことができる。

\begin{figure}[htbp]
\centering
\includegraphics[width=0.58\textwidth]{lecture_03_figure_02.png}
\caption{軸\(\hat{\boldsymbol n}\)と角度\(\theta\)が空間回転を指定する。}
\lectureframe{00:02:34}
\end{figure}

単位球面上の方向には、経度と緯度のような二つのパラメータが必要である。その方向のまわりの角度が第三のパラメータを与える。したがって回転は三パラメータの族をなす。\(2\pi\)だけ異なる角度による大域的な同一視についてはまだ気にしない。ここで重要なのは局所的なパラメータの数である。

軸と角度による記述にはいくらか冗長性がある。軸を反転させると同時に角度を反転させると同じ演算が得られ、十分に大きな角度は一周を法として還元できる。そのいずれも、恒等変換の近くでは回転が三つの連続的に変化するパラメータをもつという事実を変えない。この局所的な数は後に、行列表示における独立なパラメータの数と照合される。

群の性質を順に確認できる。
\begin{enumerate}
\item \emph{閉性。}二つの回転を続けて施すと別の回転が生じる。
\item \emph{単位元。}何もしないことが恒等回転である。
\item \emph{逆元。}\(\hat{\boldsymbol n}\)のまわりの\(\theta\)だけの回転は、同じ軸のまわりの\(-\theta\)だけの回転によって取り消される。
\item \emph{結合性。}三つの連続する回転は
\[
(R_1R_2)R_3=R_1(R_2R_3).
\]
を満たす。
\end{enumerate}

回転は可換ではない。\(x\)軸のまわりの回転に続いて\(y\)軸のまわりの回転を行うと、一般に同じ二つの演算を逆の順序で行った場合とは異なる向きが得られる。
\begin{equation}
R_x(\alpha)R_y(\beta)\ne R_y(\beta)R_x(\alpha).
\end{equation}
両方の角度が非常に小さいとき、差は二次の大きさになる。だからこそ無限小回転はほとんど可換に見える。それでも完全な群は非可換である。

非可換性は方程式を一切使わずに確認できる。本か箱を用意し、一つの面を正面として印を付け、\(x\)軸のまわりに九十度回し、次に\(y\)軸のまわりに九十度回す。元の向きに戻して順序を逆にする。印を付けた面は異なる方向を向いて終わる。小さな角度では不一致は見えにくいが、消えたわけではない。それは二つの小さな角度の積から始まる。

\begin{classroomqa}
\audiencequestion{これらの回転は同時ではなく逐次的なのか。\lecturetimestamp{00:04:45}}
\lecturerresponse{はい。群の乗法は一つの演算に続いて別の演算を行うことを意味する。それら逐次の回転の順序こそが、可換でなくなりうるものなのである。}
\end{classroomqa}

これが我々の最初の非可換群である。最終的な応用は通常の物体の色ではなく、クォークの赤、緑、青の色状態である。それは三状態の内部系であり、その興味深い変換は三つの色を混ぜ合わせる。通常の回転は、連続群が行列を通じて作用する仕方の馴染み深い模型を与える。

群それ自体は抽象的である。物理学で用いるには、それがどのような対象に作用するかを述べなければならない。{表現}は、群要素の乗法が割り当てられた行列の乗法によって再現されるように、各群要素に演算子または行列を割り当てる。異なる物理的対象は同じ群の異なる表現を担うことができる。ベクトル、スピン一の状態、スピン二分の一の状態が三つの例を与える。

\section{三次元ベクトル表現}

通常のベクトルが成分
\begin{equation}
\boldsymbol V=(V_x,V_y,V_z),\qquad V_i,\quad i=1,2,3.
\end{equation}
をもつとする。座標軸を固定したまま物理的なベクトルを回転させてもよいし、ベクトルを固定したまま座標軸を回転させてもよい。規約は異なるが、どちらの記述でも成分は\(3\times3\)行列によって混合される。任意の軸と角度の組\((\hat{\boldsymbol n},\theta)\)に対して、
\begin{equation}
V_i'=R_{ij}(\hat{\boldsymbol n},\theta)V_j,
\label{eq:l3-vector-rotation}
\end{equation}
と書く。ここで、繰り返される添字は和をとる。

繰り返される添字\(j\)はダミー添字である。式~\eqref{eq:l3-vector-rotation}は、各出力成分\(i\)に対して、明示的に
\begin{equation}
V_i'=R_{i1}V_1+R_{i2}V_2+R_{i3}V_3
\end{equation}
を意味する。自由添字\(i\)はどの成分が計算されているかを示し、繰り返される\(j\)は和をとられ、式を変えずに名前を変えてもよい。これら二つの役割を区別しておけば、長さ保存の議論における添字の誤りのほとんどを防ぐことができる。

\begin{figure}[htbp]
\centering
\includegraphics[width=0.68\textwidth]{lecture_03_figure_03.png}
\caption{回転行列は元のベクトル成分を回転後の成分に写す。}
\lectureframe{00:07:41}
\end{figure}

真の回転は長さを保存する。今のところ\(V_i\)と\(R_{ij}\)はともに実数であるから、任意のベクトルに対して
\begin{equation}
V_i'V_i'=V_iV_i
\end{equation}
が成り立つ。行列の条件を明らかにするために、同じ変換を異なるダミー添字で書くと、
\begin{equation}
V_i'=R_{ij}V_j,\qquad V_i'=R_{ik}V_k.
\end{equation}
となる。掛けて\(i\)について和をとると、
\begin{align}
V_i'V_i'
 &=R_{ij}V_jR_{ik}V_k\\
 &=R_{ij}R_{ik}V_jV_k.
\end{align}
が得られる。これが\(V\)の任意の選び方に対して\(V_jV_j\)に等しくなるのは、
\begin{equation}
R_{ij}R_{ik}=\delta_{jk}.
\label{eq:l3-orthogonality-index}
\end{equation}
のときに限る。行列記法では、
\begin{equation}
R^TR=\boldsymbol 1,
\label{eq:l3-orthogonal}
\end{equation}
であり、したがって
\begin{equation}
R^{-1}=R^T.
\end{equation}
となる。転置は単に行と列を入れ替える。

なぜ係数は\(\delta_{jk}\)でなければならないのか。クロネッカーのデルタは和の下で\(j=k\)を設定する：
\begin{equation}
\delta_{jk}V_jV_k=V_1^2+V_2^2+V_3^2.
\end{equation}
もし\(R_{ij}R_{ik}\)が\(\delta_{jk}\)と異なれば、ある方向またはある交差項\(V_jV_k\)が長さを変えてしまう。等式は一つの選ばれたベクトルに対してだけでなく、あらゆる可能なベクトルに対して要求されるから、行列の係数自体が単位行列に等しくなければならない。

式~\eqref{eq:l3-orthogonal}は、この実行列表現に対する直交性の定義条件である。行列を掛け合わせると、空間回転の乗積表が再現される。これは回転群の唯一の表現ではなく、単に最も視覚化しやすい表現にすぎない。

逆関係には直ちに幾何学的な意味がある。新しい行列方程式を解かなくても、回転を元に戻すことができる。すなわち、行と列を入れ替えればよい。\(R_1\)と\(R_2\)が二つの回転を表すなら、\(R_2R_1\)はそれらの逐次作用を表し、やはり直交している。したがって、行列の演算は、幾何学的に確立された同じ閉性、単位元、逆元、結合的な構造を保つ。

\section{スピン1とベクトル表現}

三次元表現は量子状態にも作用する。スピン1の粒子は三つの可能なスピン射影
\begin{equation}
m=+1,\quad 0,\quad -1.
\end{equation}
をもつ。したがってそのスピン状態は、確率振幅の三成分列である。粒子を回転させること、またはスピン軸を定義するために用いる磁場の向きを変えることは、回転群の三次元表現によってこれら三つの振幅を混合する。

スピン状態の数はおなじみの規則\(2s+1\)に従う。スピン0は一状態、スピン1/2は二状態、スピン1は三状態、スピン3/2は四状態をもつ。したがって状態ベクトルの次元はスピンとともに変化する。同じ物理的回転は、これらの異なる状態空間に作用するとき、異なる大きさの行列によって表現されなければならない。

スピン1の粒子に対して、基底
\begin{equation}
\lvert 1,+1\rangle,qquad
\lvert 1,0\rangle,qquad
\lvert 1,-1\rangle.
\end{equation}
を選ぶ。状態は
\begin{equation}
\lvert\psi\rangle
=c_{+}\lvert1,+1\rangle+c_0\lvert1,0\rangle+c_{-}\lvert1,-1\rangle.
\end{equation}
である。装置を回転させても第四のスピン状態が生じるのではなく、三つの係数\(c_+,c_0,c_-\)が互いに混合される。したがって対応する回転演算子は\(3\times3\)行列でなければならない。

スピン1に対して、この表現は通常の空間ベクトルと同じように変換する。これがスピン1の粒子を\emph{ベクトル粒子}と呼ぶことの正確な意味である。すなわち、その三つのスピン振幅とベクトルの三つの成分は同じ回転則に従う。量子状態を文字通りの矢印として想像する必要はなく、この記述はその成分がどのように混合されるかに関するものである。

この対応は偶然ではない。三次元ベクトルとスピン1の状態は、回転群の等価な三次元表現を担う。emph{ベクトル粒子}という言葉は、その変換性を要約したものである。量子偏極状態が粒子の内部に隠された微小な剛体の矢印であると主張しているのではない。

\section{離散対称性と連続対称性}

スピン1/2に移る前に、群の大きく分けて2つの種類を区別しておくと役に立つ。反射は単純な離散的な例を与える。恒等変換があり、反射\(P\)がある。反転するか、しないかのどちらかである。対照的に、位相
\begin{equation}
e^{i\phi}
\end{equation}
は\(\phi\)に連続的に依存し、空間回転はその軸と角度に連続的に依存する。このような連続群は、Sophus Lieにちなんで\emph{Lie群}と呼ばれ、その名は\emph{リー}と発音される。連続対称性が主な関心事となるが、離散対称性も依然として重要である。

2要素の反射群では、
\begin{equation}
P^2=1.
\end{equation}
\(1\)と\(P\)の間に連続的に調整できる反射の量は存在しない。位相群は異なる。\(e^{i\phi_1}\)と\(e^{i\phi_2}\)の間には連続的に多くの他の位相が存在する。空間回転も同じ連続的な性格を持つが、\(U(1)\)とは異なり非可換である。

スピン0は回転群の最も単純な表現を与える。スピン0の状態はあらゆる回転によって変わらないので、すべての群要素は1次元の恒等変換で表現される：
\begin{equation}
D^{(0)}(R)=1.
\end{equation}
群自体は依然として無限に多くの回転を持つ。自明なのは、スピン0の状態への表現である。

群とその表現の1つの間のこの区別は本質的である。回転群は、スカラーに作用するときはいつでも1要素の群に崩壊するわけではない。むしろ、そのすべての要素が、その特定のスカラー状態上で同じ数\(1\)によって表現されるのである。

\section{一般的なスピン1/2状態}

スピン1/2系は2つの基底状態を持つ。選ばれた\(z\)軸に対して、
\begin{equation}
\lvert\uparrow\rangle=
\begin{pmatrix}1\\0\end{pmatrix},
\qquad
\lvert\downarrow\rangle=
\begin{pmatrix}0\\1\end{pmatrix}.
\end{equation}
一般的な状態は2成分スピノルであり、
\begin{equation}
\lvert\psi\rangle=
\begin{pmatrix}\alpha_1\\\alpha_2\end{pmatrix},
\qquad
|\alpha_1|^2+|\alpha_2|^2=1.
\label{eq:l3-spinor}
\end{equation}
成分は複素数である。\((\alpha_1,\alpha_2)=(1,0)\)の選択がスピンアップであり、\((0,1)\)がスピンダウンである。一般的な規格化された組み合わせは、何か他の方向に沿って偏極したスピンを記述する。

emph{アップ}とemph{ダウン}という言葉は、軸が選ばれた後にのみ意味を持つ。ここではその軸は\(z\)である。異なる測定軸は異なる基底状態の対を持ち、ある軸に沿って確定的にアップである状態は、一般に別の基底では重ね合わせとして現れる。物理的なスピンを回転させるか、基準軸を回転させるかによって、係数はそれに応じて変化する。

二つの複素振幅には四つの実パラメータが含まれる。規格化によって一つが除かれ、物理的に無意味な全体位相によってさらに一つが除かれ、二つの実パラメータが残る。これは球面上の方向を選ぶのにちょうど必要な数である。回転を調べるために、\((\alpha_1,\alpha_2)\) とその方向の間の明示的な写像を導出する必要はない。

全体位相が除かれるのは、
\begin{equation}
\begin{pmatrix}\alpha_1\\\alpha_2\end{pmatrix}
\quad\hbox{and}\quad
e^{i\chi}
\begin{pmatrix}\alpha_1\\\alpha_2\end{pmatrix}
\end{equation}
が同じ物理的レイを表すからである。規格化は全確率を固定するのであり、通常のユークリッド空間における方向を固定するのではない。両方の同一視の後、残る2つのパラメータは、スピン偏極方向の極角と方位角と見なすことができる。

スピンが回転されると、その状態は
\begin{equation}
\alpha_i'=U_{ij}\alpha_j,
\qquad
\begin{pmatrix}\alpha_1'\\\alpha_2'\end{pmatrix}
=
\begin{pmatrix}
U_{11}&U_{12}\\
U_{21}&U_{22}
\end{pmatrix}
\begin{pmatrix}\alpha_1\\\alpha_2\end{pmatrix}.
\label{eq:l3-spinor-transform}
\end{equation}
となる。したがって、空間回転はスピノルに作用する2次元行列表現を持つ。この表現には大域的な2価性があり、それは後回しにする。局所的には、行列は以下に必要な変換性を持つ。

変換された対\((\alpha_1',\alpha_2')\)は再び可能な量子状態でなければならない。したがって、\(U\)についての最初の問いは、その明示的な軸-角度の公式ではなく、あらゆる入力スピノルの規格化を保存するという条件である。

\section{確率の保存とユニタリ行列}

回転は規格化された状態を別の規格化された状態に写さなければならない：
\begin{equation}
\alpha_i^*\alpha_i=\alpha_i'^*\alpha_i'=1.
\end{equation}
式~\eqref{eq:l3-spinor-transform}を代入すると、
\begin{align}
\alpha_i'^*\alpha_i'
 &=(U_{ij}\alpha_j)^*(U_{ik}\alpha_k)\\
 &=\alpha_j^*U_{ij}^*U_{ik}\alpha_k.
\end{align}
したがって、実ベクトルの条件の類似物は
\begin{equation}
(U^T)^*U=\boldsymbol 1.
\end{equation}
である。転置に続く複素共役はエルミート共役であり、
\begin{equation}
U^\dagger=(U^T)^*,
\end{equation}
であるから、規格化を保存する条件は
\begin{equation}
U^\dagger U=\boldsymbol 1.
\label{eq:l3-unitarity}
\end{equation}
である。

余分な複素共役が、実ベクトルの計算からの唯一の本質的な変更である。実ベクトルに対して、\(V_iV_i\)は長さである。スピノルに対して、\(\alpha_i^*\alpha_i\)は確率である。その確率を実で非負にするために複素共役が必要であり、したがって転置はエルミート共役に置き換えられる。

式~\eqref{eq:l3-unitarity}は逆行列も直ちに与える：
\begin{equation}
U^{-1}=U^\dagger.
\end{equation}
したがって、逆回転は同じ確率ノルムを保存する。通常の回転行列と同様に、ユニタリ行列の積はユニタリのままであり、表現は群の乗法を尊重する。

\begin{figure}[htbp]
\centering
\includegraphics[width=0.62\textwidth]{lecture_03_figure_04.png}
\caption{スピノル回転行列は\(U^\dagger U=\boldsymbol 1\)を満たす。}
\lectureframe{00:27:41}
\end{figure}

式~\eqref{eq:l3-unitarity}を満たす行列はユニタリである。一般の\(2\times2\)複素行列は8個の実パラメータから始まる。ユニタリ性は\(2\times2\)行列方程式であり、4個の独立な実制約を与え、4個のパラメータが残る。空間回転は3個しか持たない。\(U(2)\)全体から回転のスピノル表現を分離するには、もう1つの制限が必要である。

4個の制約は、\(U\)の2つの列を2つの複素ベクトルと見なすことによっても理解できる。各列は単位ノルムを持たなければならず、これが2個の実条件を与え、列は直交しなければならず、これが1個の複素条件すなわち2個の実条件を与える。合わせて4個である。\(U(2)\)に残る4番目のパラメータは、新しい空間回転ではなく、全体の行列式の位相であることが判明する。

\section{行列式1と\texorpdfstring{\(SU(2)\)}{SU(2)}}

追加の条件は
\begin{equation}
\det U=1.
\label{eq:l3-det-one}
\end{equation}
である。これは群の乗法と両立する。なぜなら
\begin{equation}
\det(AB)=\det A\,\det B.
\end{equation}
だからである。行列式1の2つのユニタリ行列の積もまた行列式1を持ち、単位元は行列式1を持つ。これらの行列は閉じた部分群を形成する。

逆元もまたこの部分群に属する。もし\(\det U=1\)ならば、
\begin{equation}
\det U^{-1}=\frac{1}{\det U}=1.
\end{equation}
である。したがって、行列式1を課すことは単位元を除去せず、閉性を損なわず、逆元を失わせない。それは連続パラメータを1つ減らした群を一貫して選択する。

ユニタリ性だけから
\begin{equation}
\det(U^\dagger U)=
(\det U)^*\det U
=|\det U|^2=1.
\end{equation}
が導かれる。それは\(\det U\)を\(+1\)または\(-1\)に強制しない。一般のユニタリ行列式は任意の位相であり得る：
\begin{equation}
\det U=e^{i\phi},
\qquad
\det U^\dagger=e^{-i\phi}.
\end{equation}
条件~\eqref{eq:l3-det-one}は、この他の点では自由な位相を固定し、1個の実パラメータを除去する。

\begin{classroomqa}
\audiencequestion{行列式1の条件なしで、\(U^\dagger U=\boldsymbol 1\)は\(\det U^\dagger\)が\(\det U\)に等しいこと、または\(+1\)もしくは\(-1\)のみであることを強制するか？\lecturetimestamp{00:31:21}}
\lecturerresponse{いいえ。2つの行列式は複素共役である。それらは\(e^{i\phi}\)と\(e^{-i\phi}\)であり得、その積は1である。特殊ユニタリ制限はその後、両方を1に固定する。}
\end{classroomqa}

ユニタリな\(n\times n\)行列は\(U(n)\)を形成する。\(n=1\)に対して、ユニタリ行列は単に位相\(e^{i\phi}\)であり、\(U(1)\)を与える。行列式1の部分群は\emph{特殊ユニタリ群}であり、
\begin{equation}
SU(n)=\{U\in U(n):\det U=1\}.
\end{equation}
である。2成分スピノルは\(SU(2)\)によって作用され、これは3パラメータ群である。

この記法は両方の制限を記録している：\(U\)はユニタリを意味し、\(S\)は特殊、すなわち行列式1を意味する。これらの特殊ユニタリ群は繰り返し現れる。なぜなら、それらは複素内積を保存しつつ、行列式が担う独立な共通位相を排除するからである。

局所的には、\(SU(2)\)はスピン1/2状態を回転させるのに必要な回転代数と乗法構造を同じく持つ。大域的には、先送りにされた微妙な点を思い出さなければならない：\(SU(2)\)は通常の回転群\(SO(3)\)の二重被覆であるため、\(SU(2)\)の2つの元が1つの空間回転に対応する。\editorialnote{二重被覆の定式化は、原典で言及されたが展開されなかった二値性を精密にする。}

\section{無限小回転とパウリ行列}

\begin{classroomqa}
\audiencequestion{スピン行列、すなわちパウリ行列は、この構成にどこで入ってくるのか？\lecturetimestamp{00:33:16}}
\lecturerresponse{それらは、群の元が単位元に非常に近く取られるとき、すなわち非常に小さな角度の回転に対応するときに現れる。それらはそれらの無限小回転の生成子である。}
\end{classroomqa}


\begin{equation}
U=\boldsymbol 1+\epsilon M+O(\epsilon^2),
\label{eq:l3-small-u-anti}
\end{equation}
から始める。ここで\(\epsilon\)は実で小さい。すると
\begin{equation}
U^\dagger=\boldsymbol 1+\epsilon M^\dagger+O(\epsilon^2),
\end{equation}
であり、ユニタリ性から
\begin{align}
U^\dagger U
 &=\boldsymbol 1+\epsilon(M^\dagger+M)+O(\epsilon^2).
\end{align}
が得られる。1次まででは、
\begin{equation}
M^\dagger=-M.
\end{equation}
である。したがって、式~\eqref{eq:l3-small-u-anti}の行列\(M\)は反エルミートである。

\(\epsilon\)について線形な項のみが保持されている。積
\(\epsilon^2M^\dagger M\)
は2次であるため捨てられる。これが単位元に近いところで作業する利点である：非線形な群条件が生成子に関する線形条件になる。

\(i\)の因子をくくり出す方がより便利である：
\begin{equation}
U=\boldsymbol 1+i\epsilon H+O(\epsilon^2).
\label{eq:l3-small-u-herm}
\end{equation}
複素共役は\(i\)を\(-i\)に変えるので、
\begin{equation}
U^\dagger=\boldsymbol 1-i\epsilon H^\dagger+O(\epsilon^2).
\end{equation}
である。すると
\begin{equation}
U^\dagger U
=\boldsymbol 1+i\epsilon(H-H^\dagger)+O(\epsilon^2),
\end{equation}
であり、1次のユニタリ性は
\begin{equation}
H^\dagger=H.
\end{equation}
を要求する。\(i\)の因子を露わにすると、生成子はエルミートである。

\(i\)を掛けることは許される無限小行列を変えていない；それは生成子を名付けるために用いる規約を変えただけである。物理学では通常、エルミート生成子が選ばれる。なぜなら、その固有値は実であり、角運動量のような関連する観測量がエルミートであるからである。

行列式が1であるという条件がトレースレス性を与える\editorialnote{講義中の行列式の展開では添字が一瞬乱れている。訂正された線形項は\(A_{11}+A_{22}=\operatorname{tr}A\)であり、非対角成分の積は二次の次数から始まる。}。小さな\(2\times2\)
行列\(A\)に対して、
\begin{equation}
\det(\boldsymbol 1+A)
=
\det
\begin{pmatrix}
1+A_{11}&A_{12}\\
A_{21}&1+A_{22}
\end{pmatrix}.
\end{equation}
展開すると、
\begin{align}
\det(\boldsymbol 1+A)
 &=(1+A_{11})(1+A_{22})-A_{12}A_{21}\\
 &=1+A_{11}+A_{22}+O(A^2)\\
 &=1+\operatorname{tr}A+O(A^2).
\end{align}
これを\(A=i\epsilon H\)に適用すると、
\begin{equation}
\det U=1+i\epsilon\,\operatorname{tr}H+O(\epsilon^2).
\end{equation}
したがって\(\det U=1\)は
\begin{equation}
\operatorname{tr}H=0.
\end{equation}
を要求する。トレース条件は、単位行列に比例する生成子をちょうど取り除く。一般のエルミートな\(2\times2\)行列は4つの実パラメータを持つ。すなわち2つの実対角成分と1つの複素非対角成分である。そのトレースは1つの実数の組み合わせを固定し、3つが残る。これは、微小な空間回転を行える3つの軸に対応する。

したがって\(SU(2)\)の微小生成子は、トレースレスのエルミートな\(2\times2\)行列である。3つのパウリ行列が基底をなす：
\begin{equation}
\begin{aligned}
\sigma_x&=
\begin{pmatrix}0&1\\1&0\end{pmatrix},
&
\sigma_y&=
\begin{pmatrix}0&-i\\i&0\end{pmatrix},
\\[0.4em]
\sigma_z&=
\begin{pmatrix}1&0\\0&-1\end{pmatrix}.
\end{aligned}
\end{equation}
すべての生成子は
\begin{equation}
H=a_x\sigma_x+a_y\sigma_y+a_z\sigma_z.
\end{equation}
と書ける。\(x\)軸のまわりの小さな回転は\(\sigma_x\)を選び、\(y\)軸のまわりの回転は\(\sigma_y\)を選び、その他の任意の軸のまわりの回転は対応する線形結合を選ぶ。\(\hbar=1\)とする単位系では、スピン演算子は
\begin{equation}
S_i=\frac{\sigma_i}{2}.
\end{equation}
である。

例えば、\(x\)と\(y\)の中間にある軸は、
\begin{equation}
\frac{\sigma_x+\sigma_y}{\sqrt2}.
\end{equation}
に比例する生成子を用いる。係数ベクトル\((a_x,a_y,a_z)\)の向きが回転軸を選び、その大きさが\(\epsilon\)と組み合わさって小さな角度を定める。これが、角運動量の交換関係に現れるスピン行列と有限回転演算子との間の関係である。特に、
\begin{equation}
[S_x,S_y]=iS_z
\end{equation}
およびその巡回的な相棒たちは、すでに幾何学的に見た非可換性を表している。標準的な有限回転は、微小生成子を指数関数化することで得られる\editorialnote{指数関数による有限回転の公式は、微小構成の標準的な完成である。原典は生成子の条件とそのパウリ行列基底のみを導出している。}、
\begin{equation}
U(\hat{\boldsymbol n},\theta)
=\exp\left(-\frac{i\theta}{2}\,
\hat{\boldsymbol n}\cdot\boldsymbol\sigma\right).
\end{equation}
\section{3状態の色空間}

表現の合成は後回しにしてよい。ここで空間的なスピンから離れ、クォークの色に戻る。クォークは依然としてスピン1/2をもつが、その色は独立な3状態の自由度である：
\begin{equation}
\lvert r\rangle,\qquad
\lvert g\rangle,\qquad
\lvert b\rangle.
\end{equation}
一般的な色状態は
\begin{equation}
\begin{aligned}
\lvert q\rangle_{\mathrm{color}}
={}&
\alpha_r\lvert r\rangle
+\alpha_g\lvert g\rangle
+\alpha_b\lvert b\rangle,
\\
&|\alpha_r|^2+|\alpha_g|^2+|\alpha_b|^2=1.
\end{aligned}
\end{equation}
等価に、
\begin{equation}
q=
\begin{pmatrix}
\alpha_r\\
\alpha_g\\
\alpha_b
\end{pmatrix}.
\end{equation}
二乗した大きさは、選ばれた色基底における赤、緑、青の確率である。同じ三重項表記で3つのクォーク場演算子をまとめることができる：
\begin{equation}
\psi=
\begin{pmatrix}
\psi_r\\
\psi_g\\
\psi_b
\end{pmatrix}.
\end{equation}

スピンと色は概念的に分けて保たなければならない。スピンのラベルは、その状態が物理空間の回転の下でどのように変換するかを教える。色のラベルは、実験室を回転させない内部対称性の下でどのように変換するかを教える。クォーク状態は両方のラベルを担う。ただし、現在の議論では色に集中するためにスピンを抑えている。

3つの場の演算子を1つの列に書くことは、単なる帳簿付け以上のものである。それにより、対称性を1つの行列方程式として述べることが可能になる。ちょうど空間回転が1つのベクトル成分を他の成分の線形結合に変えるように、変換は赤の場を赤、緑、青の場の線形結合に変えうる。

まず2つの色を交換する離散的な変換を想像できる。より大きな対称性は、3つすべての連続的な混合を含む。単純な交換行列は、その行列式を1に等しくするために付随する位相を必要とすることがあり、そのとき交換は特殊ユニタリ群に埋め込まれる。

2つの基底状態に対して、素朴な交換は
\begin{equation}
X=
\begin{pmatrix}0&1\\1&0\end{pmatrix},
\qquad
\det X=-1.
\end{equation}
で表される。行列全体に\(i\)を掛けると
\begin{equation}
\det(iX)=i^2\det X=1.
\end{equation}
となる。したがって、位相を首尾一貫して選べば、離散的な交換はより大きな連続群の特定の要素として現れうる。

\begin{classroomqa}
\audiencequestion{2つの基底状態を単に交換するだけの行列は、\(SU(2)\)に対して誤った行列式を持つのか。\lecturetimestamp{00:43:46}}
\lecturerresponse{素朴な交換の行列式は\(-1\)である。適切な\(i\)の因子が行列式1の代表を与えるので、交換は連続的な特殊ユニタリ変換の中に含めることができる。}
\end{classroomqa}

\section{\texorpdfstring{\(SU(3)\)}{SU(3)}と8つの色生成子}

連続的な色変換は
\begin{equation}
\begin{aligned}
q'&=Uq,
&
U^\dagger U&=\boldsymbol 1,
\\
\det U&=1,
&
U&\in SU(3).
\end{aligned}
\label{eq:l3-su3-action}
\end{equation}
である。数学はスピノルに用いたのと同じ特殊ユニタリ構成であるが、今度は\(U\)は\(3\times3\)である。これは物理空間の回転ではない。内部色空間における回転である。

ユニタリ性は色状態の規格化を保つ：
\begin{equation}
q^\dagger q
=|\alpha_r|^2+|\alpha_g|^2+|\alpha_b|^2.
\end{equation}
行列式が1であることは、連続的な混合から共通の\(U(1)\)位相方向を取り除く。得られる行列は、1つの閉じた群にとどまりながら、単純な交換、相対位相の変化、そして真に連続的な重ね合わせを行うことができる。

パラメータの数はすぐに分かる。\(3\times3\) 複素行列は9個の複素成分、すなわち18個の実パラメータをもつ。行列方程式 \(U^\dagger U=\boldsymbol 1\) は9個の実条件を課す。行列式が1であることはさらにもう1個を除く：
\begin{equation}
\dim SU(3)=18-9-1=8.
\end{equation}
したがって \(SU(3)\) は8個の独立な無限小生成子をもつ。

九つの実ユニタリティ条件は、再び\(U\)の列によって理解することができる。三つの複素列はそれぞれ単位ノルムを持ち、どの二つも直交しなければならない。これらの条件は、元の十八個の実パラメータのうち九個を取り除く。行列式の位相が十番目を取り除き、八個が残る。これは、\(8-4-1=3\)の\(SU(2)\)に対する数え方の三状態版である。

8が現れることは、最終的に8種のグルーオンを生み出すのと同じ事実である。行列式の制限が1つの方向を除いたからグルーオンは9個ではないという言い方は、単なる記憶術にすぎない。より正確には、9個の色-反色行列は、単位行列に比例する1つの色一重項方向と、8個のトレースレス方向に分解する。8個のトレースレス生成子は、QCDグルーオンが担う随伴色構造をなす。\editorialnote{一重項+八重項の定式化は、出典の意図的に発見的な行列式の注意に対する標準的で正確な版であり、そのゲージ場の構成は後回しにされる。}

この段階では、8個のグルーオンとの関係は道しるべであり、まだ完全な導出ではない。ゲージ粒子を得るには、色対称性を局所的なものに昇格させ、補償場を決定しなければならない。パラメータの数え方によってすでに示されたことは、\(SU(3)\) 色変換が単位行列から離れ始めることができる8個の独立な方向が存在することである。

\section{色対称性がQCDを制約する}

基本的な対称性の主張は、物理が式~\eqref{eq:l3-su3-action}の色変換の下で不変であるということである。これが量子色力学の中心的な内部対称性である。色成分を混合することは、形の上ではスピン上向きとスピン下向きの振幅を混合することに類似しており、それらのスピノル変換は今度は空間方向の回転に関係するが、これら3つの空間を混同してはならない。

色を対称性と呼ぶことは、色基底を変えても物理的予言が変わらないことを意味する。赤、緑、青は3つの基底ベクトルの便宜的な名前であり、3つの観測可能な色素ではない。色成分の同時 \(SU(3)\) 回転は、法則の形を変えないものでなければならない。この不変性は任意の美的特徴ではなく、QCDの組織原理である。

対称性はラグランジアンを制限する。以前の \(U(1)\) の例では、共通位相変換の下での不変性が、許される項に場とその共役の釣り合った因子を含むことを要求した。QCDでは、すべての項が同様に \(SU(3)\) 色変換の下で不変でなければならない。その要求は、相互作用、保存量、およびグルーオン場の形に帰結をもつ。

提案された項が露出した色指数をもつとしよう。式~\eqref{eq:l3-su3-action}の下でその指数は回転するので、その項はそれだけでは不変でありえない。許される項は、色指数を不変な組み合わせに縮約するか、その変化を補償する変換をもつゲージ場を用いなければならない。これが対称性原理の実際的な力である。すなわち、任意の結合を排除し、残る相互作用を組織する。

QCDの完全な構成には、さらにもう1つの段階が必要である。色変換は点ごとに変化することを許されなければならず、ゲージ場がその局所的な変化を補償しなければならない。\editorialnote{色対称性の局所化は、大域的な \(SU(3)\) の議論に対する標準的なゲージ理論的完成であり、この出典講義では展開されない。} その展開は次の段階に属する。今のところ本質的な構造は整っている。すなわち、8個の生成子をもつ特殊ユニタリ3状態対称性によって作用される色三重項である。

計算は、当面必要とされるだけ群論へと進んだ。通常空間に描ける回転から始め、ベクトルとスピノルに作用する行列へ進み、次に同じ論理を内部色空間で再利用した。次の段階では、形式論に物理的な物語全体を一度に担わせるのではなく、すでに利用可能な数学的枠組みをもって素粒子物理学へ戻ることができる。
